\subsection{反不变元}

保留前一号的记号。记 \(\varepsilon(w)\) 为元素 \(w \in W\) 的行列式。因此

\[
\varepsilon (w) = (- 1) ^ {l (w)},
\]

其中长度 \(l(w)\) 是相对于反射族 \(s_{\alpha_i}\) 取的。

定义 2. 若元素 \(x \in \mathrm{A}[\mathrm{P}]\) 满足

\[
w (x) = (- 1) ^ {l (w)}. x
\]

对所有 \(w \in W\) 成立，则称 x 在 W 下反不变。

A{[}P{]} 的反不变元素构成 A{[}P{]} 的一个 A-子模。对于任意 \(x \in A[P]\)，置

\[
J (x) = \sum_ {w \in W} \varepsilon (w). w (x).\tag{1}
\]

对于 \(x \in \mathrm{A}[\mathrm{P}]\) 和 \(w \in \mathrm{W}\)，有

\[
w (\mathrm{J} (x)) = \sum_ {v \in \mathrm{W}} \varepsilon (v). w v (x) = \varepsilon (w) \sum_ {v \in \mathrm{W}} \varepsilon (v). v (x) = \varepsilon (w). \mathrm{J} (x)
\]

且 \(J(x)\) 是反不变的。另一方面，令 \(q = \text{Card}(W)\)。对于 A{[}P{]} 的任意反不变元素 x，有 \(J(x) = q.x\)。因此，若 q 在 A 中可逆，则映射 \(q^{-1}J\) 是从 A{[}P{]} 到反不变元素子模上的投影。

设 \(\overline{\omega}_{1},\ldots,\overline{\omega}_{l}\) 为与室 C 对应的基本权。\(P\cap\overline{C}\)（分别为 \(P\cap C\)）中的元素是形如 \(n_{1}\overline{\omega}_{1}+\cdots+n_{l}\overline{\omega}_{l}\) 的权，其中 \(n_{i}\geqslant0\)（分别为 \(n_{i}>0\)），当 \(1\leqslant i\leqslant l\) 时（第 1 节，第 10 号）。另一方面，

\[
\rho = \overline {{\omega}} _ {1} + \dots + \overline {{\omega}} _ {l}
\]

是正根之和的一半（同上），所以 \(\mathrm{P} \cap \mathrm{C}\) 中的元素是形如 \(\rho + p\) 的权，其中 \(p \in \mathrm{P} \cap \overline{\mathrm{C}}\)。最后，若 \(p \in \mathrm{P} \cap \mathrm{C}\)，则 \(w(p) < p\)（对所有 \(w \neq 1\)）（\(\S 1\)，第 6 号，命题 18 的推论），因而 \(e^p\) 是 \(\mathrm{J}(e^p)\) 的唯一极大项。

命题 1. 若 2 不是 A 中的零因子，则元素 \(\mathrm{J}(e^{p})\)（其中 \(p \in P \cap C\)）构成 A{[}P{]} 的反不变元素模的一组基。

权 \(w(p)\)（其中 \(w \in \mathbf{W}\) 且 \(p \in \mathrm{P} \cap \mathrm{C}\)）两两不同。因此，\(\mathrm{J}(e^p)\)（其中 \(p \in \mathrm{P} \cap \mathrm{C}\)）线性无关。

另一方面，设 \(x = \sum_{p} x_{p} e^{p}\) 是 \(A[P]\) 的反不变元素。若 \(p_{0}\) 属于一面壁，则它在反射 \(s \in W\) 下不变，并且

\[
x = \sum_ {p} x _ {p} e ^ {p} = - s (x) = - \sum_ {p} x _ {p} e ^ {s (p)}.
\]

由此 \(2x_{p_{0}} = 0\)，故 \(x_{p_{0}} = 0\)。由于不属于任何一面壁的每个元素都能唯一写成 \(w(p)\) 的形式，其中 \(w \in W\) 且 \(p \in P \cap C\)，所以有

\[
x = \sum_ {p \in P \cap C} \sum_ {w \in W} x _ {w (p)} e ^ {w (p)}.\tag{2}
\]

由于 \(w(x) = \sum_{p} x_{p} e^{w(p)} = \varepsilon(w) \sum_{p} x_{p} e^{p}, x_{w(p)} = \varepsilon(w) x_{p}\)，从而由 (2) 得

\[
x = \sum_ {p \in P \cap C} x _ {p} J (e ^ {p}),
\]

这就完成了证明。

现在考虑元素 \(d\)，它属于代数 \(\mathrm{A}[\frac{1}{2}\mathrm{P}]\)，由下式定义：

\[
\begin{array}{l} d = \prod_ {\alpha \in R, \alpha > 0} (e ^ {\alpha / 2} - e ^ {- \alpha / 2}) \\ \quad = e ^ {\rho}. \prod_ {\alpha \in R, \alpha > 0} (1 - e ^ {- \alpha}) \\ \quad = e ^ {- \rho}. \prod_ {\alpha \in R, \alpha > 0} (e ^ {\alpha} - 1). \end{array}\tag{3}
\]

由于 \(\rho \in \mathrm{P}\)，\(d \in \mathrm{A}[\mathrm{P}]\)。

命题 2. (i) 由 (3) 定义的元素 \(d\) 是 A{[}P{]} 的反不变元素；其唯一极大项（第 2 号，定义 1）是 \(e^{\rho}\)，且 \(d = \mathrm{J}(e^{\rho})\)。

\begin{enumerate}
\def\labelenumi{(\roman{enumi})}
\setcounter{enumi}{1}
\item
  对任意 \(p \in \mathbb{P}\)，元素 \(\mathrm{J}(e^p)\) 被 \(d\) 唯一地整除，且商 \(\mathrm{J}(e^p)/d\) 是 \(\mathrm{A}[\mathrm{P}]\) 中在 \(\mathrm{W}\) 下不变的元素。
\item
  若 2 不是 A 中的零因子，则乘以 d 是从 A{[}P{]} 中 W 不变元素集到 A{[}P{]} 中反不变元素集的双射。
\end{enumerate}

我们知道，对于 \(1 \leqslant i \leqslant l\)，反射 \(s_i = s_{\alpha_i}\) 保持除 \(\alpha_i\) 外的正根集合不变，且 \(s_i(\alpha_i) = -\alpha_i\)（第 1 节，第 6 号，命题 17 的推论 1）。因此，

\[
\begin{array}{l} s _ {i} (d) = (e ^ {- \alpha_ {i} / 2} - e ^ {\alpha_ {i} / 2}). \prod_ {\alpha \in R, \alpha > 0, \alpha \neq \alpha_ {i}} (e ^ {\alpha / 2} - e ^ {- \alpha / 2}) \\ = - d = \varepsilon (s _ {i}). d. \end{array}
\]

由于 \(s_{i}\) 生成 W，这证明了 (i) 的第一个断言。(i) 的第二个断言立即由 (3) 和引理 2 得出，因为 \(1 - e^{-\alpha}\) 的唯一极大项是 1（对 \(\alpha \in R\) 且 \(\alpha > 0\)）。

现在假设 A = Z。根据命题 1，

\[
d = \prod_ {p \in \mathrm{P} \cap \mathrm{C}} c _ {p} \mathrm{J} (e ^ {p}) \text {with} c _ {p} \in \mathbf {Z}.\tag{4}
\]

另一方面，显然

\[
d = e ^ {\rho} + \sum_ {q <   \rho} c _ {q} ^ {\prime} e ^ {q}.\tag{5}
\]

若 \(p \in \mathrm{P} \cap \mathrm{C}\) 且 \(p \neq \rho\)，则 \(p > \rho\)，并且根据 (5)，\(e^p\) 在 \(d\) 中的系数为零。因此 \(c_p = 0\)。此外，比较 (4) 和 (5) 中 \(e^\rho\) 的系数可得 \(c_\rho = 1\)，从而 \(d = \mathrm{J}(e^\rho)\)。

继续假设 A = Z。设 \(p \in P\)、\(\alpha \in R\)，并令 M 为相对于子群 \(\{1, s_{\alpha}\}\) 的 W 的右陪集的一组代表元。于是

\[
\mathrm{J} (e ^ {p}) = \sum_ {w \in \mathrm{M}} \varepsilon (w) e ^ {w (p)} + \sum_ {w \in \mathrm{M}} \varepsilon (s _ {\alpha} w) e ^ {s _ {\alpha} w (p)}.
\]

现在 \(s_{\alpha}w(p) = w(p) - \langle \alpha^{\prime},w(p)\rangle \alpha = w(p) + n_w\alpha\)，其中 \(n_w\in \mathbf{Z}\)。因此

\[
J (e ^ {p}) = \sum_ {w \in M} \varepsilon (w) e ^ {w (p)} \left(1 - e ^ {n _ {w} \alpha}\right).
\]

若 \(n_w \geqslant 0\)，显然 \(1 - e^{n_w\alpha}\) 被 \(1 - e^\alpha\) 整除；当 \(n_w < 0\) 时这仍然成立，因为 \(1 - e^{n_w\alpha} = -e^{n_w\alpha}(1 - e^{-n_w\alpha})\)。因此，\(J(e^p)\) 被 \(1 - e^\alpha\) 整除（在 \(\mathbf{Z}[P]\) 中）。

根据引理 1，\(\mathbf{Z}[\mathrm{P}]\) 是唯一分解环，且元素 \(1 - e^{\alpha}\)（其中 \(\alpha \in \mathbb{R}\) 且 \(\alpha > 0\)）两两互素。因此，\(\mathrm{J}(e^{p})\) 在 \(\mathbf{Z}[\mathrm{P}]\) 中被乘积 \(\prod_{\alpha > 0}(1 - e^{\alpha})\) 整除，因而也被 \(d = e^{-\rho}\prod_{\alpha > 0}(1 - e^{\alpha})\) 整除。

回到一般情形，通过从 Z 到 A 的标量扩张，由上面得到 \(d = \mathrm{J}(e^{\rho})\)，且每个元素 \(\mathrm{J}(e^{p})\) 都被 d 整除。~由于 \(e^{\rho}\) 是 d 的唯一极大项，第 2 号的备注表明，存在唯一的元素 \(y \in A[P]\)，使得 \(\mathrm{J}(e^{p}) = dy\)，并且立即得到 y 在 W 下不变，从而 d 与 \(\mathrm{J}(e^{p})\) 都是反不变的。这证明了 (i) 和 (ii)。

最后，若 2 不是 A 中的零因子，则第 2 号备注和命题 1 蕴含 (iii)。

备注。1）若 2 不是 A 中的零因子，容易验证 d 是 A{[}P{]} 中以 \(e^{\rho}\) 为其极大项的唯一反不变元素。

\begin{enumerate}
\def\labelenumi{\arabic{enumi})}
\setcounter{enumi}{1}
\tightlist
\item
  第 2 号引理 2 表明，商 \(\mathrm{J}(e^p) / d\)（当 \(p\in \mathrm{P}\cap \mathrm{C}\) 时）的唯一极大项是 \(e^{p - \rho}\)。
\end{enumerate}

